\chapter{Data Augmentation}
\label{chap:augmentation}

The first augmentation attempt (shift and mirror) never beat the plain baseline. Rather than
stop there, five more methods were tried, each tested alone first, then the winners combined
--- and one of them turned out to be the best single model in the entire project. This chapter
covers all six methods, why each one was chosen, and the full, honest result of testing them.

\section{Shift and mirror: the first attempt}

Only seizure (class 1) windows are created, and only from the real, filtered signal of the
training recordings --- never from recordings held out for testing, and never anything
synthesized from nothing.

\begin{itemize}[leftmargin=1.4em]
  \item \textbf{Shift.} Extra 2-second windows starting 0.25, 0.5, and 0.75 seconds after each
    regular window's start time (the regular window grid has a 1-second hop), kept under the
    same rule used everywhere else: at least 50\% seizure coverage.
  \item \textbf{Mirror.} Every seizure window, regular or shifted, with its left and right
    electrodes swapped: Fp1$\leftrightarrow$Fp2, F7$\leftrightarrow$F8, F3$\leftrightarrow$F4,
    T3$\leftrightarrow$T4, C3$\leftrightarrow$C4, T5$\leftrightarrow$T6, P3$\leftrightarrow$P4,
    O1$\leftrightarrow$O2. The three midline channels, Fz, Cz, and Pz, are left unchanged.
\end{itemize}

Across the 21 recordings this adds 8{,}230 synthetic seizure windows on top of the 1{,}244
original ones: 3{,}493 shifted, 1{,}244 mirrored, and 3{,}493 both shifted and mirrored. These
were chosen first because both stay entirely inside real, recorded signal --- nothing is
synthesized, so nothing physiologically implausible can be fabricated by the transformation
itself. Mirroring is anatomically justified specifically for absence seizures (generalized,
bilaterally synchronous discharges, so a left-right-swapped window is still a valid seizure
pattern); shifting is consistent with a tolerance already built into the official labeling rule.

The result, reported honestly: shift+mirror raised recall for every algorithm (SVM: 86 to 90 of
93 found) but also raised false alarms sharply (SVM: 16 to 32) and lowered event F1 (SVM: 0.894
to 0.862). It never beat the plain, unbalanced baseline for any of the four algorithms.

\section{Five more methods, tried the same way}

A negative result for one method is not a verdict on augmentation in general. Five more
methods were built, each respecting the same rule as shift and mirror --- training folds only,
built from real signal, one new synthetic window per real seizure window (no extra time
offsets, so each method adds exactly 1{,}244 windows, directly comparable to one another):

\begin{itemize}[leftmargin=1.4em]
  \item \textbf{Amplitude scaling.} Multiply the whole window by a random gain between 0.6 and
    1.6, simulating the kind of electrode-to-electrode or session-to-session gain variation a
    real recording can have.
  \item \textbf{Noise.} Add Gaussian noise to the window, scaled to that window's own signal
    level, simulating ordinary sensor noise.
  \item \textbf{Time masking.} Blank out a short random contiguous span inside the window (up to
    25\% of it) to zero, forcing the model not to rely on any one instant.
  \item \textbf{Channel dropout.} Zero out 1 to 3 random EEG channels for the whole window,
    simulating electrodes that failed or were never connected.
  \item \textbf{Combined.} The windows from whichever of the above turned out to help, merged
    into one dataset --- built only after the individual results were in, not before.
\end{itemize}

The broader menu of EEG augmentation methods also includes frequency-domain transforms (Fourier
surrogates, frequency shift), mixing methods (mixup, cutmix, pasting segments), and generative
models (GANs, VAEs, diffusion). These were set aside for the same reason shift and mirror were
originally preferred over them: generative methods need far more than 93 seizures to avoid
memorizing or fabricating unrealistic signal, and mixing methods risk blending two unrelated
seizure segments into a composite that does not correspond to any real physiological event. The
five methods above were chosen because each one models a specific, plausible real-world
imperfection (gain drift, sensor noise, a dead channel, a brief dropout) rather than an
arbitrary transformation.

\section{What actually helped}

\begin{table}[H]
  \centering
  \small
  \caption{Every augmentation method, every algorithm, event F1 at its own best threshold.
    Bold marks the best result for that algorithm; a dagger marks a result that beats the
    algorithm's own unbalanced baseline.}
  \label{tab:augmentation-full}
  \begin{tabular}{llrrr}
    \toprule
    Algorithm & Method & Event F1 & Found/93 & False alarms \\
    \midrule
    SVM & unbalanced (baseline) & 0.894 & 86 & 16 \\
    SVM & shift + mirror & 0.862 & 90 & 32 \\
    SVM & amplitude & 0.882 & 88 & 23 \\
    SVM & noise & 0.889 & 87 & 20 \\
    SVM & time masking & 0.870 & 87 & 25 \\
    SVM & \textbf{channel dropout} & \textbf{0.901}$^\dagger$ & 88 & 18 \\
    SVM & combined (noise + dropout) & 0.889 & 87 & 20 \\
    \midrule
    Random Forest & unbalanced (baseline) & 0.867 & 87 & 26 \\
    Random Forest & shift + mirror & 0.854 & 84 & 26 \\
    Random Forest & amplitude & 0.867 & 87 & 24 \\
    Random Forest & noise & 0.872$^\dagger$ & 86 & 23 \\
    Random Forest & time masking & 0.859 & 79 & 15 \\
    Random Forest & \textbf{channel dropout} & \textbf{0.880}$^\dagger$ & 82 & 14 \\
    Random Forest & combined (noise + dropout) & 0.878$^\dagger$ & 81 & 13 \\
    \midrule
    KNN & unbalanced (baseline) & 0.828 & 86 & 37 \\
    KNN & \textbf{noise} & \textbf{0.832}$^\dagger$ & 81 & 29 \\
    KNN & combined (noise + dropout) & 0.829$^\dagger$ & 81 & 30 \\
    \midrule
    Decision Tree & unbalanced (baseline) & 0.290 & 91 & 561 \\
    Decision Tree & \textbf{combined (noise + dropout)} & \textbf{0.296}$^\dagger$ & 93 & 619 \\
    Decision Tree & noise & 0.295$^\dagger$ & 93 & 570 \\
    \bottomrule
  \end{tabular}
\end{table}

Two methods actually work: \textbf{channel dropout} and \textbf{noise}. Channel dropout gives
the single best result in this entire project --- SVM with channel-dropout augmentation reaches
event F1 \textbf{0.901} (88 of 93 found, 18 false alarms, threshold 0.90), beating the plain SVM
baseline (0.894) by the same margin the much more expensive SVM+CNN ensemble achieved
(\Cref{chap:ensemble}), except this is one plain model, not a blend of two. Channel dropout also
helps Random Forest (0.867 to 0.880). Noise helps Random Forest, KNN, and (marginally)
Decision Tree.

Time masking hurt every algorithm it was tried on, and amplitude scaling was mixed (helped SVM,
did nothing for Random Forest). Combining the two winners, noise and channel dropout, into one
dataset did \emph{not} beat channel dropout alone for either SVM (0.889 vs.\ 0.901) or Random
Forest (0.878 vs.\ 0.880) --- the single best method outperformed the combination, a useful
reminder that more augmentation is not automatically better augmentation.

\textbf{Why channel dropout specifically works}: \Cref{chap:preprocessing} showed that no single
EEG channel dominates a trained model's decisions (the heaviest channel, O1, carries under 10\%
of total importance), consistent with absence seizures being a generalized, whole-scalp pattern.
Channel dropout trains the model on windows missing 1--3 channels, which cannot destroy the
seizure signature --- plenty of real signal remains in the other 16--18 channels. Time masking,
by contrast, removes a slice of \emph{every} channel at the same instant, which can erase the
one thing that actually identifies an absence seizure: its rhythm over time. Losing space is
survivable for this problem; losing time is not.

\begin{summarybox}
Six augmentation methods were tried: the original shift+mirror, four new methods tested in
isolation (amplitude scaling, noise, time masking, channel dropout), and one combined dataset
built from the two methods that actually helped. Only channel dropout and noise improved
results; channel dropout gives the best result in the whole project, SVM at event F1 0.901,
beating the plain SVM baseline (0.894) as a single model, no blending required. Combining the
two winners did not beat the best one alone. Time masking consistently hurt, and the original
shift+mirror method never beat the baseline for any algorithm --- a real negative result,
reported alongside the real positive one.
\end{summarybox}
